\section{HAETAE 구조 및 오류주입 연구 동향}
\label{sec:background}

\subsection{HAETAE 전자서명}

\code{HAETAE}는 $\Rq=\mathbb{Z}_q[x]/(x^{256}+1)$ 위에서 동작하는
격자 기반 전자서명이다~\cite{haetae}. 본 연구의 분석 대상인
\code{HAETAE}-120에서 $q=64513$이며, 비밀 성분 $\svec_1$은 세 개의
비밀 다항식으로 구성되어 총 768개 계수를 포함한다. 구현에서는 공개
성분을 포함하는 길이 4의 벡터를 이용하므로 응답 $z_1$은 네 개 다항식,
즉 1024개 계수로 표현된다.

서명자는 seed와 counter로부터 초구 난수 $(\yvec_1,\yvec_2)$와
bimodal 부호 $b$를 생성한다. 이후 공개행렬과 난수로부터 챌린지 $c$를
유도하고, $c\svec_1$과 $c\svec_2$를 계산하여 식
~\eqref{eq:intro-response}의 응답을 생성한다. 생성된 응답이 서명 경계
$B_1$과 bimodal 경계 $B_0$를 만족하지 않으면 해당 시도를 거부하고
새로운 난수로 서명을 다시 수행한다.

거부 샘플링은 서명 응답의 분포가 비밀키에 의존하지 않도록 만드는
절차이다. 따라서 응답 계산뿐만 아니라 거부 여부를 결정하는 프로그램
흐름도 오류주입 공격의 대상이 될 수 있다. 또한 서명 단계에서 사용하는
경계 $B_0$와 $B_1$은 검증 단계의 경계 $B_2$와 목적이 다르므로,
서명 과정의 거부 판정이 생략되었을 때 표준 검증만으로 모든 비정상 응답을
식별할 수 있다고 단정하기 어렵다.

\subsection{격자 서명에 대한 오류주입 공격}

Espitau 등은 Fiat--Shamir with Aborts 계열 전자서명의 거부 루프를
조기에 종료시켜 비밀정보에 의존하는 서명을 획득하는 공격을 제안하였다
~\cite{espitau2016}. Groot Bruinderink와 Pessl은 결정론적 격자 서명을
대상으로 클럭 글리치에 의해 생성된 정상 및 결함 서명의 차분을 이용하였다
~\cite{bruinderink2018}. Ravi 등은 덧셈 명령어 생략으로 nonce와 비밀
성분 사이의 관계가 노출되는 공격을 분석하였으며~\cite{ravi2019determinism},
이후 NTT 상수 변조 등 다양한 오류 지점으로 연구 범위가 확대되었다
~\cite{ravi2023ntt}.

ML-DSA에 대해서는 정정 오류와 명령어 생략을 이용한 비밀정보 복원
연구가 수행되었다~\cite{krahmer2024,elghamrawy2023}. Jendral 등은
hedged ML-DSA 환경에서 단일 결함 서명으로부터 비밀정보를 추론하는
조건을 분석하였다~\cite{jendral2024}. 이러한 연구는 공개 챌린지와
오류 응답 사이에 남는 대수적 관계가 격자 서명 구현의 주요 공격 표면임을
보인다.

\code{HAETAE}에 대해서는 Lee 등~\cite{lee2026haetae}이 LSB 추출,
공개행렬 언패킹, bimodal 부호 및 샘플링 seed와 관련된 네 개의 오류
지점을 분석하였다. 해당 연구는 응답을 생성하기 이전의 연산과 거부
샘플링 구조에서 발생할 수 있는 오류를 제시하고, 부분 이중 연산과 서명
후 검증 등의 대응방안을 논의하였다. Shin 등은 challenge sampling
절차의 공개계수와 오류 응답 사이의 관계를 이용하여 ML-DSA와
\code{HAETAE}에서 단일 결함 서명으로 위조에 필요한 비밀정보를
복구하는 공격을 제시하였다~\cite{shin2026public}. 본 연구는 이와
달리 응답식의 두 항과 거부 판정을 중심으로 취약 지점을 분류하고,
T1의 구현 의존적 물리 실현성을 실험한다.

\subsection{기존 연구의 한계}

기존 연구에서 사용하는 소프트웨어 오류 에뮬레이션은 목표 연산이
생략되었을 때의 논리적 결과를 결정론적으로 확인하는 데 적합하다. 하지만
소스 코드에서 하나의 항을 제거하는 것과 실제 프로세서에서 한 번의
글리치로 동일한 벡터 연산 전체를 제거하는 것은 서로 다른 문제이다.
컴파일된 명령어의 배치, 반복문의 크기, load--compute--store 순서 및
버퍼의 이전 값에 따라 물리 오류가 남기는 중간값이 달라질 수 있기
때문이다.

또한 연구용 펌웨어는 오류 효과를 확인하기 위해 챌린지, 내부 응답 또는
참 비밀값을 별도의 디버그 채널로 출력하는 경우가 많다. 이러한 계측은
오류 지점과 정보 노출 사이의 인과관계를 검증하는 데 유효하지만, 표준
인터페이스만을 사용하는 종단 공격과는 검증 범위가 다르다. 본 연구는
SW 에뮬레이션을 취약 조건과 대수적 영향의 검증에 사용하고, T1에
대해서만 물리 실현성을 추가로 평가한다. 직렬화 서명만을 이용한 공격과
위조 검증은 현재 결과의 필수 전제가 아니라 후속 공격 고도화 범위로
구분한다.
